\begin{corollary}[序列层折叠的熵不降：每步保留 \texorpdfstring{$1$}{1} 比特自由度]\label{cor:Phi_m-entropy-no-drop}
对任意 \(m\ge 2\)，令 \(Z_m\subset E_m^{\ZZ}\) 为图 \(G_m\) 的双边边移位、\(\mathcal{L}_m:Z_m\to X_m^{\ZZ}\) 为逐边取标记映射（定义 \ref{def:G_m}），并记 \(Y_m=\mathcal{L}_m(Z_m)\)（定理 \ref{thm:Phi_m-sofic-graph}）。
则：
\begin{enumerate}
  \item \(Z_m\) 与二元全移位 \(\{0,1\}^{\ZZ}\) 共轭，因而 \(h_{\mathrm{top}}(Z_m)=\log 2\)；
  \item 标记图 \(G_m\) 右解析（同一顶点出发的两条出边标记不同；定理 \ref{thm:pom-right-resolving}），从而 \(\mathcal{L}_m\) 是有限对一因子码（每个输出序列的 lift 数至多为 \(|V_m|=2^{m-1}\)）；
  \item 因而拓扑熵保持：
  \[
  \boxed{\ h_{\mathrm{top}}(Y_m)=h_{\mathrm{top}}(Z_m)=\log 2\ }.
  \]
\end{enumerate}
该结论给出“窗口层多对一压缩”与“序列层记忆补偿”的严格命题：\(\Phi_m\) 将窗口层丢失的纤维自由度迁移到时间方向的有限状态记忆中，而不在熵意义下压缩每步 \(1\) 比特的自由度。
\leanverified{LabeledPath}
\leanverified{labeledPath\_initial\_injective}
\leanverified{labeledPath\_card\_le\_state\_card}
\leanverified{labeledPath\_count\_le\_states}
\leanverified{paper\_labeled\_path\_lift\_bound}
\leanverified{paper\_Phi\_m\_entropy\_no\_drop}
\end{corollary}
\begin{proof}
\textbf{(1)}--\textbf{(2)} 见定理 \ref{thm:Phi_m-sofic-graph}（共轭）与右解析性；\textbf{(3)} 有限对一因子码熵保持可参见 \cite{LindMarcus1995,Kitchens1998SymbolicDynamics}。
\end{proof}

\begin{remark}[时间记忆作为隐寄存器]\label{rem:Phi-m-memory-implicit-register}
窗口层的多对一压缩并不等价于时间层的信息损失：当 $m\ge 3$ 时，定理 \ref{thm:Phi_m-conjugacy-threshold} 给出有限窗逆码 $\Psi_m$，从而被折叠抹去的自由度可迁移到时间方向的有限状态记忆并在有限延迟后恢复。这提供了“第三寄存器”的另一种实现方式：其不必以静态轴显式存储，而可作为同步滤波器的内部状态隐式承载。在规范群胚口径下，该隐寄存器对应局部截面之间的粘接数据（见小节 \ref{subsec:pom-gauge-groupoid}）。
\end{remark}
